% >>> typage-template ctexart@2026.10.1 variant=xiaosi
\documentclass[UTF8,fontset=fandol,zihao=-4,a4paper]{ctexart}
\usepackage[margin=2.5cm]{geometry}
\usepackage{gbt7714}
\bibliographystyle{gbt7714-numerical}
% <<< typage-template
\usepackage{amsmath,amssymb}
\usepackage{graphicx}
\usepackage{booktabs}
\usepackage[hidelinks]{hyperref}

\title{论文标题}
\author{第一作者\thanks{某大学某学院，\texttt{first.author@example.com}} \and 第二作者\thanks{某研究所}}
\date{\today}

\begin{document}
\maketitle

\begin{abstract}
用几句话说明研究问题、方法和主要结论。
\end{abstract}

\section{引言}
介绍研究背景与问题，列出本文的贡献。中文排版的基本用法可参考文献\cite{liu2013latex}，
\LaTeX{} 的通用约定见文献\cite{lamport1994latex}。

\section{方法}
给需要引用的公式编号：
\begin{equation}
  \label{eq:objective}
  \min_{\theta} \; \frac{1}{n} \sum_{i=1}^{n} \ell\bigl(f_\theta(x_i), y_i\bigr) + \lambda \lVert \theta \rVert_2^2 .
\end{equation}
式~\eqref{eq:objective} 是训练目标。

\section{实验结果}
表~\ref{tab:results} 汇总了实验结果。

\begin{table}[htbp]
  \centering
  \caption{示例数据集上的结果}
  \label{tab:results}
  \begin{tabular}{lcc}
    \toprule
    方法     & 准确率（\%） & 用时（秒） \\
    \midrule
    基线方法 & 81.2         & 12.4       \\
    本文方法 & 84.7         & 10.9       \\
    \bottomrule
  \end{tabular}
\end{table}

\section{结论}
总结全文并展望后续工作。

\bibliography{refs}
\end{document}
